\scanpage{19}
\uline{Définition de $\beta$}\qquad On définit
\[
\beta(\Pi,\widetilde\Pi,p,\mathcal G,a)=(\mathcal E,\mathcal G,\pi,Z)\quad\text{ainsi :}
\]
$\mathcal E=$ groupe des couples $(g,\bar g)$, où $g\in\mathcal G$, $\bar g\in\operatorname{Aut}(\widetilde\Pi)$ compatible avec $g$
\[
\pi=\operatorname{Aut}(\widetilde\Pi/\Pi)
\]
$\mathcal E\mathbin{\unsure{-|\longrightarrow}}\mathcal G$ par $(g,\bar g)\longmapsto g$, d'où
\diagram{\begin{tikzcd}[ampersand replacement=\&]
1 \arrow[r,dash] \& \pi \arrow[r] \& \mathcal E \arrow[r,dash] \& \mathcal G \arrow[r] \& 1
\end{tikzcd}}
$Z=\operatorname{Ob}\widetilde\Pi$, avec opération évidente de $\mathcal E$ sur $Z\ldots$

Variantes profinies\ldots

\textcircled{8} Considérons, avec notations \unsure{du} \S 7,
\[
\left(\diagram{\begin{tikzcd}[ampersand replacement=\&]
1 \arrow[r,dash] \& \pi \arrow[r,dash] \& \mathcal E \arrow[r] \& \mathcal G \arrow[r,dash] \& 1
\end{tikzcd}}
\quad\text{et }Z\quad\mathcal E\text{-ens}\ne\varnothing\text{ sur lequel }\pi\text{ opère librement}\right)
\]
$\longleftrightarrow(\Pi,\widetilde\Pi,p,\mathcal G,a)$. On va interpréter la donnée d'un triple $(\inserted{\chi,}\tau,T,\Phi)$, où $\tau$ est un $\mathcal E$-ens, $T$ un groupe $(\simeq\widehat{\mathbf Z}\;?)$, \inserted{$\chi:\mathcal G\longrightarrow\operatorname{Aut}(T)$ \unsure{$(\simeq\widehat{\mathbf Z}^{*}\;\text{canon.}?)$}} et
\diagram{\begin{tikzcd}[ampersand replacement=\&]
\Phi:\tau \arrow[r] \& \operatorname{Hom}_{\mathrm{gr}}(T,\pi)
\end{tikzcd}}
compatible avec op. de $\mathcal E$ (opérant
\mspage{18}
